\section{Metodologia}\label{sec:metodologia}

Esta seção apresenta a metodologia aplicada na investigação, sobre como como LLMs podem aprender, ou como se diferenciam nesta representação, da tarefa de solução de labirintos em ASCII, por meio da análise das mudanças nas ativações internas das redes neurais, após diferentes estratégias de \textit{fine-tuning}. O objetivo desta seção é detalhar os procedimentos adotados, de modo a garantir a replicabilidade do estudo, incluindo preparação dos dados, modelos utilizados, métricas utilizadas e estratégias de análise estatísticas.

\subsection{Dados}

Foram utilizadas amostras do conjunto de dados de treino, disponibilizado por \cite{dao2025alphamaze}, representados em ASCII, uma vez que esses dados estão disponíveis publicamente e permitiram uma redução considerável no tempo que seria empregado na elaboração dos datasets e ajuste dos modelos.

\begin{itemize}
    \item \textbf{Formato geral dos exemplos}: cada exemplo é um prompt textual que (i) instrui o modelo sobre a tarefa; (ii) apresenta a grade completa do labirinto com origem e destino; e (iii) define que a saída deve conter apenas os tokens de movimento.

    \item \textbf{Tokenização do labirinto}: a descrição é composta por tokens delimitados no padrão \verb!<|...|>!, incluindo:
    \begin{itemize}
        \item Coordenadas: \verb!<|row-col|>! (ex.: \verb!<|0-0|>!, \verb!<|3-2|>!) indicam a posição absoluta da célula na grade.
        \item Paredes: um único token resume as paredes presentes na célula, combinando direções com underscore, por exemplo \verb!<|up_wall|>!, \verb!<|left_right_wall|>!, \verb!<|up_left_right_wall|>!.
        \item Marcadores especiais: \verb!<|origin|>! (origem) e \verb!<|target|>! (destino).
        \item Separador: \verb!<|blank|>! atua como espaçamento explícito entre grupos de tokens/células, reduzindo ambiguidade de fronteiras.
        \item Vocabulário de movimentos (saída do modelo): \verb!<|up|>!, \verb!<|down|>!, \verb!<|left|>!, \verb!<|right|>!.
    \end{itemize}

    \item \textbf{Padrão por célula}: cada célula segue tipicamente \verb!<|row-col|><|wall_pattern|><|blank|>!, com este último podendo incluir \verb!<|origin|>! ou \verb!<|target|>! quando aplicável.

    \item \textbf{Semântica das paredes}: o token de paredes compacta o estado de fronteiras da célula. Por exemplo, \verb!<|up_left_wall|>! indica a existência de paredes acima e à esquerda; \verb!<|left_right_wall|>! bloqueia as laterais; a ausência de uma direção implica passagem livre.
    
    \item \textbf{Formato da resposta}: o modelo deve emitir apenas a sequência de tokens de movimento, separados por espaços, que leva da origem ao destino. Em alguns prompts incentiva-se ``\textit{pensar passo a passo}'', mas a saída final deve conter exclusivamente os tokens de movimento.
\end{itemize}

\subsection{Modelos e Estratégias de Fine-tuning}

Serão selecionados LLMs abertos com até 8 bilhões de parâmetros, compatíveis com a biblioteca LLM-MRI. Existe também a possibilidade de utilizar os mesmos modelos empregados no AlphaMaze, visto que esses modelos estão disponíveis publicamente, o que facilitaria a comparação de resultados. O processo de especialização dos modelos será realizado por meio de diferentes estratégias de \textit{fine-tuning}, conforme discutido na literatura:
\begin{itemize}
    \item \textbf{Supervised Fine-Tuning (SFT)}: ajuste supervisionado utilizando exemplos de labirintos e suas soluções.
    \item \textbf{Group Relative Policy Optimization (GRPO)}: ajuste baseado em preferências coletivas, visando aprimorar o raciocínio e a autocorreção dos modelos.
\end{itemize}
Cada modelo será treinado separadamente em cada estratégia, utilizando os mesmos dados de entrada para garantir comparabilidade.

\subsection{Avaliação de Desempenho}

O desempenho dos modelos será avaliado antes e após o \textit{fine-tuning}, utilizando métricas objetivas como taxa de sucesso na resolução dos labirintos, número médio de passos até a solução e precisão na geração dos comandos. Para garantir rigor na avaliação, os resultados serão comparados com abordagens tradicionais e com benchmarks da literatura, como o MazeBench.

\subsection{Análise das Ativações Neuronais}

A análise das ativações internas será conduzida com o auxílio da biblioteca LLM-MRI \cite{costa2024llmmri}. Para cada modelo e estratégia de ajuste, serão coletados os vetores de ativação das camadas intermediárias ao processar diferentes labirintos. Uma redução de dimensionalidade será aplicada para visualizar e comparar as distribuições das ativações antes e depois do treinamento. O objetivo será identificar padrões emergentes, alterações estruturais e possíveis mecanismos de raciocínio espacial desenvolvidos pelos modelos.

\subsection{Hipóteses do Estudo}

Este estudo parte das seguintes hipóteses principais:
\begin{itemize}
    \item[H1:] Estratégias distintas de \textit{fine-tuning}, como SFT e GRPO, produzirão padrões diferentes de ativação interna, refletindo abordagens variadas de raciocínio espacial e representação da tarefa.
    \item[H2:] A análise das ativações neuronais, por meio da LLM-MRI, permitirá identificar alterações estruturais e padrões emergentes associados ao aprendizado da tarefa, contribuindo para a compreensão dos mecanismos internos dos modelos.
\end{itemize}

\subsection{Discussão e Interpretação dos Resultados}

Os resultados quantitativos e qualitativos serão analisados de forma integrada, buscando compreender como as diferentes estratégias de \textit{fine-tuning} influenciam tanto o desempenho quanto as representações internas dos LLMs. A discussão considerará limitações, possíveis vieses e implicações para o desenvolvimento de modelos mais interpretáveis e eficientes em tarefas que exigem raciocínio espacial.
